% Reviewed standalone English write-up of the final construction.
% UTF-8; no external bibliography file is required.
% Suggested build: latexmk -pdf -outdir=pdf maximum_weight_sparse_subgraph_en_reviewed.tex
% Mathematical review and reproducible finite checks: October 11, 2026.
\documentclass[11pt,a4paper]{article}
\usepackage[T1]{fontenc}
\usepackage[utf8]{inputenc}
\usepackage[english]{babel}
\usepackage{lmodern}
\usepackage[margin=26mm]{geometry}
\usepackage{amsmath,amssymb,amsthm,mathtools}
\usepackage{microtype}
\usepackage{enumitem}
\usepackage{float}
\usepackage{needspace}
\usepackage[hidelinks,unicode]{hyperref}
\usepackage{bookmark}

\newtheorem{theorem}{Theorem}[section]
\newtheorem{lemma}[theorem]{Lemma}
\newtheorem{proposition}[theorem]{Proposition}
\newtheorem{corollary}[theorem]{Corollary}
\theoremstyle{definition}
\newtheorem{definition}[theorem]{Definition}
\theoremstyle{remark}
\newtheorem{remark}[theorem]{Remark}
\newcommand{\clj}[2]{\operatorname{cl}_{M_{#1}}(#2)}
\newcommand{\Bj}[2]{B_{#1}(#2)}
\newcommand{\rank}{\operatorname{rk}}
\newcommand{\Sparse}{\textsc{Sparse}}
\newcommand{\Filter}{\textsc{Filter}}
\newcommand{\Circuits}{\textsc{Circuits}}
\newcommand{\FlowZero}{\textsc{FlowZero}}
\newcommand{\softO}{\widetilde O}
\newcommand{\indeg}{d^-}
\newcommand{\birth}{b}
\newcommand{\botval}{\bot}
\setlist{itemsep=3pt,topsep=5pt}
\setlength{\emergencystretch}{2em}
\floatstyle{ruled}
\newfloat{algorithm}{htbp}{loa}
\floatname{algorithm}{Algorithm}
\hypersetup{
 pdftitle={A Deterministic Almost-Linear-Time Algorithm for the Maximum-Weight (k,l)-Sparse Subgraph Problem},
 pdfauthor={}
}
\title{A Deterministic Almost-Linear-Time Algorithm\\
 for the Maximum-Weight $(k,\ell)$-Sparse Subgraph Problem\\[4pt]
 \large The full classical range and a uniform additive input bound}
\author{}
\date{October 11, 2026}

\begin{document}
\maketitle

\begin{abstract}
For each fixed $k\geq1$ and $0\leq\ell<2k$, we give a single
deterministic word-RAM algorithm with worst-case running time
$O(m+n^{1+o(1)})$ for the maximum-weight $(k,\ell)$-sparse
subgraph problem, assuming that the edges are already ordered
by nonincreasing weight. The algorithm does not depend on an
exponent-error parameter, and the coefficient of the $m$ term
does not depend on the recursion depth.

In the upper range $k<\ell<2k$, we reduce the number of mandatory
tight bodies geometrically using a fixed constant number of local
greedy layers. Samples are generated implicitly from shared ordered
incidence lists. Edges spanned by earlier selected edges of their
body pair are deleted permanently. A separate saturation argument
bounds priority-queue head repairs by $O(Dp+qD+f)$ at a stage with
$q$ bodies, $p$ physical vertices, and $f$ known auxiliary constraints.
Fixed first bindings and activation times preserve chronological
correctness across stages without repeated ancestor queries.
The union of all samples is an $O(n)$-edge certificate.

This English version organizes the final construction and all
required preliminary lemmas into a standalone presentation.
The generic natural representation is used only to prove safety;
the algorithm performs deterministic combinatorial operations.
\end{abstract}

\tableofcontents

\section{Problem, model, and main theorem}
\label{sec:problem}

Let $G=(V,E)$ be a loopless undirected multigraph with $n$ vertices
and $m$ edges. Vertex identifiers belong to $\{1,\ldots,n\}$.
The input supplies a strict total edge order
\[
 e_1\succ e_2\succ\cdots\succ e_m
\]
consistent with nonincreasing weight; input order breaks ties.
Nonpositive-weight edges may be omitted without changing the
maximum attainable weight. We therefore work with positive weights.
All subsequent greedy-output and certificate equalities refer to this
retained positive-weight input. An empty input is returned immediately;
we may therefore assume $n\geq2$.

\begin{definition}
An edge set $I$ is \emph{$(k,j)$-sparse} if
\[
 |I[X]|\leq\max\{k|X|-j,0\}
 \qquad(X\subseteq V).
\]
A set $X$ with $|X|\geq2$ is \emph{tight} if
$|I[X]|=k|X|-j$. A \emph{block} is an inclusion-wise maximal
nontrivial tight vertex set.
\end{definition}

For $0\leq j<2k$, these independent edge sets form a matroid,
denoted by $M_j$. Write $\Bj jA$ for the ordered matroid-greedy
output on an edge set $A$, and $A_{>e}$ for the edges of $A$
preceding $e$. The greedy output is a maximum-weight independent
set and is determined by the supplied order.

\begin{theorem}\label{thm:main}
For fixed $k\geq1$ and $0\leq\ell<2k$, the maximum-weight
$(k,\ell)$-sparse subgraph problem on ordered edges has a single
deterministic word-RAM algorithm with worst-case running time
\[
 \boxed{O(m+n^{1+o(1)})}.
\]
The program and the coefficient of the $m$ term depend only on
$k,\ell$. No choice of algorithm indexed by an arbitrarily small
positive exponent is required.
\end{theorem}

The word length is $\Omega(\log(n+m+2))$, with a sufficiently large
constant to address the polynomial-size tables used below.
Vertex and edge identifiers, root-prefix timestamps, and the
polynomially bounded integers in the combinatorial data structures
fit in a machine word. The constants and the $o(1)$ term may
depend on fixed $k,\ell$. We do not claim the same bound when
$k$ is part of the input, or a corresponding almost-linear space
bound. In particular, constant-time initializable direct-address
tables are allowed \cite{initializable}.

The flow subroutine in \cite{mcf} uses fixed-point numbers with
$\log^{O(1)}z$ bits on an instance encoded with $z$ bits
(its Section~2, ``Computational Model''). On our sparse networks
$z=N\log^{O(1)}N$. Simulating those arithmetic operations on
ordinary machine words has a polylogarithmic overhead, which is
absorbed in $N^{1+o(1)}$. This simulation is confined to the sparse
subroutines; the operations charged directly to the original $m$
records use only a constant number of machine-word operations.

The proof proceeds in three parts. First, we construct a
deterministic almost-linear optimizer for sparse inputs and a
compact history of all tight prefix blocks. Second, we establish
a local sampling and body-compression lemma. Third, we execute
sampling implicitly and compose the stages with fixed endpoint
bindings. The complete algorithm appears in
Section~\ref{sec:complete}.

\subsection{Algorithmic primitives}

The following published algorithms are used as subroutines.
All their parameters other than the instance size are fixed here.
\begin{enumerate}
 \item Exact minimum-cost flow and circulation with polynomially
 bounded integer capacities and costs can be computed
 deterministically in almost-linear time \cite{mcf}.
 \item An edge set known to decompose into $k$ forests can be
 decomposed deterministically in almost-linear time \cite{forest}.
 \item A rooted unit-capacity digraph known to have rooted
 connectivity at least $a$ admits $a$ actual edge-disjoint spanning
 out-arborescences, computable in $\softO(N+A)$ time for fixed
 $a$ \cite{arborescences}. The bound reported in
 \cite{arbo-study} is
 $O(Aa\log N+Na^4\log^2N)$ for exact rooted connectivity $a$.
 If only a lower bound is given, attach a new root by $a$ parallel
 arcs to the old root, apply that algorithm, and remove these arcs.
 \item For a fixed positive minimum-cut value, the required
 inclusion-wise minimal directed minimum-cut sets can be extracted
 in $\softO(N+A)$ time from the minset-poset representation
 \cite{minset,minset-2026}. For value zero, SCCs suffice.
\end{enumerate}
Here $N$ and $A$ denote network vertices and arcs. Each explicit
sparse instance is relabelled densely on $\{1,\ldots,N\}$; its
separate metadata retain the original vertex identities and timestamps.
The reductions below use these algorithms only on sparse networks.
The orientation and rooted-connectivity perspective also appears
in the recognition algorithms of \cite{recognition}.

\section{Tight sets and consecutive sparsity matroids}
\label{sec:matroids}

\begin{lemma}[Uncrossing and closure]\label{lem:uncross}
Let $I$ be independent in $M_j$. If tight sets $A,B$ have at
least two common vertices, then $A\cap B$ and $A\cup B$ are tight,
and there is no edge of $I$ between $A\setminus B$ and
$B\setminus A$. Distinct blocks therefore intersect in at most
one vertex and have disjoint internal edge sets.
For a new edge $e=uv\notin I$,
\[
 e\in\clj jI
 \quad\Longleftrightarrow\quad
 u,v\text{ belong to a common tight set}.
\]
If $j<k$, distinct blocks are vertex-disjoint.
Every nontrivial tight set has no isolated internal vertex.
\end{lemma}
\begin{proof}
Supermodularity of induced edge counts and sparsity give
\[
 k|A|+k|B|-2j
 \leq |I[A\cup B]|+|I[A\cap B]|
 \leq k|A\cup B|+k|A\cap B|-2j.
\]
The two ends are equal, so all inequalities are equalities.
For closure, a violated sparsity inequality in $I+uv$ must
contain both endpoints and must have been tight in $I$.
Conversely, adding an internal edge to a tight set violates its
bound. With a one-vertex intersection, two tight sets would
violate the union bound when $j<k$.
For the last assertion, a tight two-vertex set has
$2k-j>0$ edges. In a larger tight set, removing an isolated
vertex would leave the same edge count but lower the sparsity
bound by $k$, a contradiction.
\end{proof}

\begin{lemma}[Spanned tight graphs]\label{lem:spanned-tight}
Let $I$ be independent in $M_j$, and let $T$ be an independent
tight graph on $X$, where $|X|\geq2$. If every edge of $T$ belongs
to $\clj jI$, then $X$ is contained in a block of $I$.
\end{lemma}
\begin{proof}
Let $\mathcal Z$ be the blocks of $I$ meeting $X$ in at least two
vertices, and put $Y=X\cup\bigcup_{Z\in\mathcal Z}Z$.
Every edge of $T$ outside these blocks belongs to $I$, by
Lemma~\ref{lem:uncross}. If $\mathcal Z$ is empty, this gives
$T\subseteq I$, so $X$ itself is tight in $I$.
Otherwise the internal edge sets of the blocks are disjoint, and
\begin{align*}
 k|X|-j=|T|
 &\leq |I[Y]|-\sum_{Z\in\mathcal Z}|I[Z]|
       +\sum_{Z\in\mathcal Z}(k|X\cap Z|-j)\\
 &\leq k|Y|-j-\sum_{Z\in\mathcal Z}(k|Z|-j)
       +\sum_{Z\in\mathcal Z}(k|X\cap Z|-j)\\
 &\leq k|X|-j.
\end{align*}
The last inequality uses
$|Y|\leq |X|+\sum_Z(|Z|-|X\cap Z|)$.
Equality throughout forces $|I[Y]|=k|Y|-j$.
Thus $Y$ is tight and a block containing $Y$ contains $X$.
\end{proof}

\begin{lemma}[Disjoint circuits]\label{lem:circuits}
Let $1\leq j<2k$. If $I$ is independent in $M_{j-1}$, its $M_j$-circuits are
pairwise edge-disjoint. Their supports are precisely the
inclusion-wise minimal nontrivial $(k,j-1)$-tight sets $W$, and
each such circuit is
\[
 I[W],\qquad |I[W]|=k|W|-(j-1).
\]
Consequently, $\Bj jI$ is obtained by deleting the last edge of
each circuit in the supplied order.
\end{lemma}
\begin{proof}
Every edge of a circuit must be internal to a violated set;
otherwise its internal edges give a proper dependent subset.
Independence in $M_{j-1}$ and dependence in $M_j$ determine the
displayed edge count. If two circuits shared an edge, their
supports would meet in two vertices. Uncrossing would make their
intersection weakly tight and hence dependent in $M_j$,
contradicting minimality unless the circuits were equal.
Every dependence contains a circuit, and disjoint circuits
interfere with none of the other greedy deletions.
\end{proof}

\begin{lemma}[Greedy compatibility]\label{lem:compatibility}
For every ordered $E$ and $1\leq j<2k$,
\[
 \Bj j{\Bj{j-1}E}\subseteq\Bj jE.
\]
\end{lemma}
\begin{proof}
Suppose the stronger greedy algorithm accepts $e$ when restricted
to the weaker greedy basis. Let $D$ be the weaker accepted prefix
before $e$; thus $D+e$ is weakly independent. Let $\mathcal Z$
be the maximal weakly tight blocks of $D$.
Every earlier edge rejected by the weaker algorithm is internal
to a member of $\mathcal Z$. The stronger restricted greedy
algorithm accepts every edge of $D$ outside these blocks:
otherwise Lemma~\ref{lem:circuits} would give a weakly tight
circuit support containing such an edge. The edge $e$ is also
outside the blocks.

Let $A$ be the stronger full greedy prefix before $e$. Construct
$K$ from the edges of $D+e$ outside the blocks and the sets
$A[Z]$, $Z\in\mathcal Z$. Then $A+e\subseteq K$.
For $|X|\geq2$, put
\[
 \mathcal Z_X=\{Z\in\mathcal Z:|Z\cap X|\geq2\},
 \qquad Y=X\cup\bigcup_{Z\in\mathcal Z_X}Z.
\]
If $\mathcal Z_X$ is nonempty, weak sparsity of $D+e$ bounds
the outside-block contribution by
\[
 k|Y|-(j-1)
 -\sum_{Z\in\mathcal Z_X}(k|Z|-(j-1)).
\]
The inside contribution is at most
$\sum_{Z\in\mathcal Z_X}(k|X\cap Z|-j)$. Using
$|Y|\leq |X|+\sum_{Z\in\mathcal Z_X}(|Z|-|X\cap Z|)$, their sum is at most
$k|X|-(j-1)-|\mathcal Z_X|\leq k|X|-j$.
If $\mathcal Z_X$ is empty, $K[X]$ is a subset of the stronger
restricted accepted prefix together with $e$, and is independent.
Thus $K$ is independent in $M_j$, so $A+e$ is independent as well.
\end{proof}

\section{Sparse optimization and the complete prefix history}
\label{sec:sparse}

This section supplies the static algorithms used later on
explicit graphs of $O(N)$ or $O(N\log(N+2))$ edges. These graphs
contain only registered physical vertices and constant-size
body frames; the original dense input is never passed to them.

\subsection{Rooted-cut reductions}

Let $I$ be $(k,t)$-sparse. If $t\geq k$, decompose $I$ into
$k$ forests $T_1,\ldots,T_k$. A tight set $C$ satisfies
\[
 \sum_{i=1}^k (|C|-1-|T_i[C]|)=t-k.
\]
Every summand is nonnegative, so at least one of the first
$t-k+1$ forests is connected on $C$.
If $0<t<k$, decompose $I$ into $t$ forests and $k-t$
pseudoforests. Tightness implies connectivity of $C$ in
every forest part, so the first forest suffices.
This mixed decomposition can also be constructed almost linearly:
orient $I$ with indegree at most $k$, add $k-\indeg(v)$ root
arcs to each vertex, and extract $t$ edge-disjoint spanning
out-arborescences. Removing their root arcs gives the forest
parts; the remaining orientation has indegree at most $k-t$
and can be split into that many pseudoforests.

Centroid-decompose the selected forest components. For every
connected tight $C$, some call has a vertex set $U$ and centroid
$u$ with $u\in C\subseteq U$: take the first centroid in $C$.
Over all calls, vertex-set sizes sum to $O(N\log(N+2))$.
A fixed bounded-indegree orientation permits induced edges
to be listed with the same bound by scanning at most $k$
incoming edges per vertex and using membership marks.

For $t\geq k$, root every local forest at $u$, remove $u$, and
add $k-\indeg(v)$ arcs $r\to v$, retaining the indegrees
measured before removing $u$. The resulting digraph $R$ satisfies
\begin{equation}\label{eq:upper-root}
 \rho_R(X)=k|X|-|I[X\cup\{u\}]|\geq t-k=:\delta.
\end{equation}
Nonempty $\delta$-cut sets correspond exactly to tight sets
containing $u$.
For $0<t<k$, orient locally with capacities
$b(u)=k-t$ and $b(v)=k$ otherwise, and add root arcs to fill
these capacities. Then
\begin{equation}\label{eq:lower-root}
 \rho_R(X)=k|X|-t\mathbf1_{u\in X}-|I[X]|\geq0.
\end{equation}
Every nonempty zero-cut contains $u$ and is a tight set.
Indeed, a root-avoiding set not containing $u$ has cut value at
least $t>0$, while the singleton $\{u\}$ has cut value $k-t>0$.
The required orientation follows from capacitated Hall
inequalities and is found by a sparse flow reduction.
For $t=0$, use one global bounded-indegree orientation and
root completion; its cut formula is
$\rho_R(X)=k|X|-|I[X]|$.

\subsection{Listing the circuits}

To compute $\Circuits(j,I)$, apply the reductions with $t=j-1$.
In each call, enumerate the inclusion-wise minimal nonempty
root-avoiding cut sets of the target value (zero or $\delta$),
and translate them back to vertex sets. If there are no such
sets, the call contributes no candidate.
For positive $\delta$, the rooted reduction to the global
minset-poset problem adds $\delta$ arcs from every nonroot
vertex to $r$, a new vertex $z$, and $\delta$ arcs in each
direction between $r$ and $z$. The global minimum-cut value
is $\delta$, and cuts on the original nonroot vertices are unchanged.
Every cut containing $r$ receives either the $\delta$ arcs from $z$
or the $\delta$ arcs of a missing original vertex; a cut containing
$z$ but not $r$ receives the $\delta$ arcs $r\to z$.
Thus no smaller cut is introduced. Keep only minimal cuts avoiding
both $r$ and $z$. For value zero, use source SCCs having no entering
arc from outside the SCC; the root SCC is excluded.

Minimal root-avoiding minimum-cut sets are disjoint: a nonempty
intersection would, by cut submodularity, itself be a
minimum-cut set. Candidate sizes therefore sum to $O(|U|)$
per call, even when the centroid is added to every candidate.
Every true circuit support appears in its covering centroid
call. Process candidates by increasing vertex-set size,
accepting a candidate precisely when none of its internal
edges has already been marked; mark those edges on acceptance.
A nonminimal candidate contains a smaller true circuit.
A candidate sharing an edge with an accepted circuit contains
that circuit, by uncrossing and minimality. This procedure
therefore lists exactly the circuits of
Lemma~\ref{lem:circuits}, in $N^{1+o(1)}$ time.

\subsection{All tight prefix blocks}

Fix an independent ordered set $F=\{f_1,\ldots,f_a\}$, and let
$F_t=\{f_1,\ldots,f_t\}$. If $C$ is tight in a prefix, it is
tight in the final set and $F[C]=F_t[C]$: a later internal
edge would violate final independence. Its birth time is
\[
 \birth(C)=\max\{i:f_i\in F[C]\}.
\]
Hence the forest and centroid decompositions of final $F$
cover every prefix tight set.

Build each local rooted network from final $F$ and delete
original edges in reverse order. An arc deletion $a\to b$
is implemented as
\begin{equation}\label{eq:transfer}
 a\to b\ \longmapsto\ r\to b.
\end{equation}
This preserves the indegree total and can only increase
root-cut values. In the upper reduction, deleting an edge
from the removed centroid simply adds a root arc.

For zero cuts, the maximal zero-cut is the set unreachable
from $r$. For reachability, keeping the old arc and adding
the new root arc is equivalent to replacement.
Each vertex becomes reachable at most once, giving linear
total scanning work in a local call.

For $\delta>0$, extract $\delta$ spanning out-arborescences
$P_1,\ldots,P_\delta$. After root arcs are removed, let
$Q_1,\ldots,Q_\delta$ be their forests, and let $Z$ contain
the other arcs. A nonempty $X$ is a $\delta$-cut exactly when
each $Q_i[X]$ is connected and no $Z$-arc enters $X$:
each arborescence contributes one entering arc per induced
forest component, and all additional contributions are
nonnegative.

Maintain a partition of live vertices, refining it until
every part is connected in every $Q_i$. If an extra arc
enters a part from outside, delete its head from the live
set and propagate further refinements. Every possible
$\delta$-cut remains live and contained in one part:
connectivity forces containment through refinements, and
an entering extra arc excludes its head from every candidate
cut contained in that part. At stability, the parts themselves
are $\delta$-cuts, hence are exactly the maximal ones.

If a deleted arc belongs to $P_i$, its subtree is reattached
to the root. The arborescence remains valid, while $Q_i$
loses one edge. If the deleted arc is extra, its replacement
root arc excludes its head. Thus only splitting and deletion
are needed throughout the reverse history.

Dynamic forests and smaller-side relabeling give a
$\softO(N+A)$ total bound for each local history. Here is an
implementation of the refinement and propagation work, including
arcs that become external when a part splits.

Maintain each $Q_i$ as a decremental forest, cutting also every
forest edge whose endpoints acquire different part labels and
every edge incident to a deleted vertex. These additional cuts
do not remove an edge internal to any surviving candidate cut:
such a candidate is contained in one current part and uses only
live vertices. Consequently every maintained forest component
is contained in one part. Keep an explicit component identifier
for each live vertex, retaining the identifier of the larger
component on a cut and relabelling the smaller component.
Deterministic balanced Euler-tour trees give component sizes,
cuts, and enumeration of the smaller component in
$O(\log(N+2)+\text{enumerated size})$ time.
Each vertex changes its component identifier only
$O(\log(N+2))$ times in each of the fixed number of forests.

For every part and every forest, store a balanced dictionary of
the component identifiers occurring in that part, with their
vertex counts. If two components occur, enumerate the smaller
of two such components and move it to a new part. It has at most
half the old part's vertices. Update the dictionaries and scan
the original forest and extra-arc incidences of moved vertices.
Cut forest edges made external by the split, and queue the head
of each extra arc newly entering a part. Scan all extra arcs once
initially. When a queued head is deleted, cut its forest
incidences, remove it from the dictionaries, and scan its extra
incidences; in particular an outgoing extra arc may now exclude
its live head. Duplicate queue entries are checked against the
current labels and live flags before they are used. Perform all
cross-part forest cuts before choosing the next refinement.

Part labels, like forest-component labels, move only to smaller
classes, so a live vertex changes them $O(\log(N+2))$ times.
Each original incidence is therefore scanned
$O(\log(N+2))$ times, each forest edge is cut at most once,
and each vertex is deleted at most once. Dictionary and forest
operations add at most polylogarithmic factors. This proves
the stated $\softO(N+A)$ total bound without charging a scan of
a surviving large part after every update.

Record stable parts in a splitting/deletion tree, suppressing
unchanged parts and storing vertex membership by the tree rather
than copying every part. Splits have linear total tree size;
unary shrinkages are charged to the vertices deleted at those
shrinkages. Thus each local history has linear size.

To compile the local trees globally, assign each original edge
to its smallest local tight ancestor, if one exists. This ancestor
is found from the endpoint leaves by LCA preprocessing of the
membership tree; in the upper reduction an edge incident to the
centroid uses the other endpoint's leaf. The centroid is common
to all the local tight sets and is not duplicated as a membership
leaf. Ancestors that do not represent stable tight sets are
skipped. Internal edge sets are laminar within a local tree,
so births are computed by subtree maxima of assigned edge indices.
These assignments and computations take $\softO(N+A)$ time.
Form a timestamped auxiliary graph on original
edge-elements and local tight nodes: connect each edge to
its assigned node at that node's birth, and each node to
its first tight parent at the parent's birth.

At time $t$, retain only active components containing at least
one activated tight node. Their projections onto original
edge-elements are exactly the internal edge sets of the
maximal blocks of $F_t$. Indeed, a shared internal edge permits
tight-set uncrossing and joins only one global block.
Conversely, every global block is a maximal local candidate
in a centroid call covering it, so all its internal edges
are connected there.
An isolated edge leaf, or an edge group with no activated tight
node, is not a block. This distinction is necessary, for example,
when a single accepted edge is not yet a tight two-vertex graph.

Process equal timestamps in batches and record the resulting
union history. Suppress auxiliary nodes and unary nodes that
do not change an edge group. Preserve the first activation
of a one-edge tight block separately from its edge leaf.
There are $O(a)$ genuine block nodes and merges.
A nontrivial tight set has no isolated internal vertex, so
its vertex set is determined by its internal edge endpoints.
This yields the compact history in $N^{1+o(1)}$ time.

\subsection{A constant-time chronological closure index}

Orient final $F$ with indegree at most $k$. For a historical
tight block $C$, define its ports by
\[
 \Pi(C)=\{v\in C:\indeg_{F[C]}(v)<k\}.
\]
Since the internal indegree deficit sums to $j$,
$|\Pi(C)|\leq j$. A nonport vertex is \emph{saturated}.
It can be saturated in only one current maximal block,
because those blocks have disjoint internal edge sets.

Compute port lists by merging internal indegree maps up the
edge-group hierarchy, always moving the map of the smaller
edge group, measured by its number of edge leaves. Each map has
at most twice that many vertex entries, and the edge-group size
at least doubles whenever it is moved. This gives
$O(N\log(N+2))$ total entry visits even when maps share keys.
Store the deficient entries separately, so an entire large
map is never rescanned at every ancestor. Update this list only
at touched keys; at a genuine tight node it has at most $j$
entries. Deterministic dictionaries give total cost
$\softO(N)$.
For a vertex of final indegree $k$, let its anchor be the
lowest genuine block node containing all its incoming edge
leaves. The anchor is found by a constant number of LCA
queries in preprocessing, followed by the nearest genuine-block
ancestor lookup. If the incoming leaves have no common genuine
block ancestor, the vertex has no anchor; an artificial root
joining different history trees is never a block. At time $t$,
its highest activated ancestor is
the saturated owner $o_t(v)$, if one exists.
Explicit smaller-class relabeling on this known union tree
costs $O(N\log(N+2))$ in total and permits $O(1)$ owner queries;
the specialized union-find of \cite{static-union} is another
implementation.

For each distinct pair of ports, precompute
\[
 \beta(u,v)=\min\{\birth(C):u,v\in\Pi(C)\},
\]
with value $+\infty$ when the pair is absent. There are
$O(N)$ such pairs. Two stable counting sorts by endpoint
identifier, followed by merging with the query-pair list,
give every offline $\beta$ answer in $O(N+Q)$ time.
For a new edge $uv\notin F_t$, closure holds exactly when
\begin{enumerate}
 \item $o_t(u)=o_t(v)\ne\botval$;
 \item one endpoint has a saturated owner and the other is
 a port of that owner; or
 \item $\beta(u,v)\leq t$.
\end{enumerate}
These are precisely the cases in which both endpoints lie
in a common current block. Each query uses constant work.
For $j=0$, closure is membership of both endpoints in the
global unreachable set, whose entry times were recorded above.

Consequently, for fixed independent $F$, the operation
\[
 \Filter(j,F,E)=
 \{e\in E:e\notin\clj j{F_{>e}}\}
\]
costs $O(N^{1+o(1)}+|E|)$ deterministically.
The block history and port lists themselves are computed
without scanning a separate dense edge set.

\subsection{The sparse optimizer}

\begin{proposition}\label{prop:sparse}
For fixed $k\geq1$, $0\leq j<2k$, and $|E|=O(N\log(N+2))$, the ordered greedy
output $\Bj jE$ can be computed deterministically in
$N^{1+o(1)}$ time. Mandatory independent edges may be placed
first in the order.
\end{proposition}
\begin{proof}
For $j=0$, independence is equivalent to assigning each edge
to one endpoint, with at most $k$ assignments per vertex.
Build a circulation network with an edge-node $x_i$ for
each ordered edge $e_i=uv$, arcs $s\to x_i$ of capacity one
and cost $-(a-i+1)$, zero-cost unit arcs $x_i\to u,x_i\to v$,
arcs $v\to t$ of capacity $k$, and a zero-cost return arc
$t\to s$ of capacity $kN$, where $a=|E|$.
An integral minimum-cost circulation selects the maximum
weight independent set for these strictly ordered rank costs.
Matroid greedy depends only on the order, so its output is
exactly $\Bj0E$ for the original weights as well.
The distinct positive surrogate weights make the optimal edge
set unique; alternative optimal routings of its endpoint
assignments do not affect the selected set.
All numbers are polynomially bounded, and \cite{mcf} gives
the required almost-linear running time.

For $j>0$, repeatedly perform
\[
 \begin{aligned}
 B&\gets\Sparse(j-1,E),\\
 F&\gets B\setminus
 \{\text{last edge of each }C\in\Circuits(j,B)\},\\
 E&\gets\Filter(j,F,E),
 \end{aligned}
\]
until $E=F$. Initially take the previous round's set to be empty.
The filter retains $F$, because $F$ is independent.
The filter never deletes an edge of the full greedy optimum
$O$: a deleted edge is spanned by earlier input edges.
By Lemmas~\ref{lem:circuits} and \ref{lem:compatibility},
every produced $F$ is contained in $O$.

Let $F_0$ be the previous round's set. Deleting input edges
cannot turn a previously accepted weaker greedy edge into
a rejected edge, so $F_0\subseteq B$. Since $F_0\subseteq O$,
the stronger restricted greedy algorithm also accepts its
edges, and $F_0\subseteq F$.
Every stronger circuit in $B$ contains at least two edges
outside $F_0$. Its last edge cannot belong to $F_0$.
If that last edge were its only edge outside $F_0$, the
previous filter would already have deleted it.
The circuits are edge-disjoint, so at most
$|B\setminus F_0|/2$ edges are deleted. With $r=|O|$,
\[
 |F|-|F_0|\geq
 \frac{|B|-|F_0|}{2}\geq\frac{r-|F_0|}{2}.
\]
The missing optimum size halves every nonfinal round.
There are $O(\log(N+2))$ rounds. Once $F=O$, the next
filter leaves exactly $F$.

The recurrence is
$T_0(N)=N^{1+o(1)}$ and
$T_j(N)\leq O(\log(N+2))(T_{j-1}(N)+N^{1+o(1)})$.
Its depth is fixed, so $T_j(N)=N^{1+o(1)}$.
Independent mandatory edges placed first are all accepted.
\end{proof}

\section{The lower range: an ordered union certificate}
\label{sec:lower}

For $0\leq\ell\leq k$, the sparsity matroid is the union of
$\ell$ graphic and $k-\ell$ bicircular matroids; the latter
have pseudoforests as their independent sets
\cite{recognition}. We use the following general certificate.

\begin{lemma}\label{lem:union-certificate}
Let $U=N_1\vee\cdots\vee N_k$ be a union of ordered matroids
on $E$, and let $R$ be an integer with
$R\geq\max\{1,\sum_i r_i(E)\}$.
Set $L=1+\lceil\log_2(R+1)\rceil$.
Create $L$ cycles of layers, each cycle containing
$N_1,\ldots,N_k$ in that order. Insert each edge into the
first layer that accepts it, or discard it if all reject.
The union $H$ of the layers contains $B_U(E)$, and
$|H|\leq RL$.
\end{lemma}
\begin{proof}
Every earlier layer spans every element stored in a later
layer in the earlier layer's matroid, at every input prefix.
Suppose $U$-greedy accepts $e$ but every layer rejects it.
Put $A=E_{>e}$ and choose a minimizer $X\subseteq A$ of
the union rank formula
\[
 r_U(A)=|A\setminus X|+\sum_i r_i(X).
\]
Write $Q_i=N_i/X$. Comparing $X$ with $X\cup Y$ shows that
every $Y\subseteq A\setminus X$ satisfies
$|Y|\leq\sum_i r_{Q_i}(Y)$.
Thus $A\setminus X$ is independent in the union of the
quotients and has size at most $R$.
Moreover, $e$ is not a loop in at least one quotient;
otherwise choosing $X+e$ in the rank formula for $A+e$
would show that $e$ does not increase union rank.

All layer sets in the rest of this argument are evaluated just
before $e$. Let $K_c$ be the stored elements outside $X$ in cycles
strictly after cycle $c$. Each layer $B_{c,i}\setminus X$
spans $K_c$ in $Q_i$. Hence
\[
 |K_c|\leq\sum_i r_{Q_i}(K_c)
 \leq\sum_i|B_{c,i}\setminus X|
 =|K_{c-1}|-|K_c|.
\]
It follows that $|K_c|\leq |K_{c-1}|/2$.
Since $|K_0|\leq R$, we have $|K_{L-1}|<1$.
Every layer in the final cycle is therefore contained
in $X$, yet spans $e$. This makes $e$ a loop in every
$Q_i$, a contradiction.
Thus no greedy optimum edge is discarded. Removing
greedy-rejected edges does not change the greedy output,
and each cycle has size at most $R$.
\end{proof}

\begin{proposition}\label{prop:lower}
The lower range has a deterministic
$O(m+n^{1+o(1)})$ algorithm.
\end{proposition}
\begin{proof}
Use Lemma~\ref{lem:union-certificate} with $R=kn$.
Within each layer, store explicit component identifiers
and component lists, relabelling the smaller list on a merge.
A graphic layer rejects $uv$ exactly when its endpoints
are in the same component. A bicircular layer rejects
$uv$ exactly when both endpoints belong to unicyclic
components, which may be the same component.
Tree-component vertices are marked when their component
becomes unicyclic; every vertex receives this mark at most
once in a layer.
Closure queries therefore take worst-case constant time,
and maintenance costs $O(n\log(n+2))$ per layer.

Within each of the two layer types, earlier layers span
later layers, so rejecting layers form an initial segment.
First test only the final layer of each type that is present.
If every present type rejects, discard the edge in constant time.
Otherwise binary-search the accepting types to find the
first accepting layer, and choose the earliest among them.
Binary searches occur only for retained edges, of which
there are $O(n\log(n+2))$.
The total construction cost is
$O(m+n\log^2(n+2))$.
Apply Proposition~\ref{prop:sparse} to the resulting
$O(n\log(n+2))$-edge certificate.
\end{proof}

\section{Bodies, physical pins, and logical block size}
\label{sec:bodies}

For the remainder of the construction, assume
$k<\ell<2k$ and set $c=2k-\ell\geq1$.
A \emph{body} is a mandatory tight graph. Its \emph{pins}
are physical vertices, whose identities are retained even
when the internal body graph is replaced by a small frame.
Different bodies at the start of a stage are vertex-disjoint.
Their mandatory bases have union $P_0$.

\subsection{Frames and the natural representation}

The natural representation of the sparsity matroid
\cite{natural} can be described using a generic matrix
$U$ with $k$ rows $U_v$ per physical vertex and $\ell$ columns.
An edge row is supported on its endpoints and satisfies
$a_{e,u}U_u+a_{e,v}U_v=0$.
The nullspace of a tight body's internal basis is exactly
\begin{equation}\label{eq:motions}
 \xi_v=U_v\alpha_A\quad(v\in A),
 \qquad \alpha_A\in\mathbb R^\ell.
\end{equation}
The stacked matrix $U$ on a nontrivial body has column rank
$\ell$, and its independent tight basis has $k|A|-\ell$ rows.
Its nullspace therefore has dimension $\ell$ and equals the
displayed $\ell$-dimensional space, which is already annihilated
by every internal row. In particular the body parameter is unique.
All generic statements can be made simultaneously on a
sufficiently large universe, padded with isolated vertices
if necessary. No such matrix is constructed by the algorithm.

A constant-size tight frame can be built explicitly.
Let $a=\max\{2,\lceil\ell/c\rceil\}$.
On $a$ vertices, the complete multigraph with $c$ copies
per pair contains a tight graph with $ka-\ell$ edges.
For proper subsets of size $x\geq2$,
\[
 kx-\ell-c\binom{x}{2}
 =-\frac c2(x-2)(x-\ell/c)\geq0.
\]
Retain any $ka-\ell$ edges. Enlarge the frame by adding
vertices with $k$ incident edges, at most $c$ per old
vertex pair; $ca\geq k$ guarantees sufficient capacity.
This preserves sparsity and tightness.
Choose the final constant frame size at least
\begin{equation}\label{eq:frame-width}
 h=\max\{k,2,\lceil\ell/c\rceil\}.
\end{equation}

Attach a fresh pin $v$ to $k$ distinct frame vertices.
The enlarged graph is tight, so these attachments impose
exactly $\xi_v=U_v\alpha_A$. Pins need only be registered
when an actually retained edge uses them. Before registration,
their private attachments are understood implicitly.
Registering a fresh pin creates no new maximal block
apart from extending the block containing its frame.
To see this, let $C$ be a maximal tight block containing the
fresh pin before any additional incident edge is inserted.
A nontrivial tight set has no isolated internal vertex, so $C$
contains at least one of the pin's frame neighbors. It therefore
meets the enlarged tight frame in at least two vertices.
Uncrossing and maximality force $C$ to contain the whole enlarged
frame. Smaller tight subsets need not contain all the attachments:
for $(k,\ell)=(2,3)$ a single attachment is already tight on its
two endpoints.

\begin{lemma}[Exact disjoint-body count]\label{lem:body-count}
For an edge set $J$ joining distinct bodies, let $p_i(S)$
be the number of distinct pins in body $A_i$ incident to
$S\subseteq J$. Then $J$ is independent after contracting
$P_0$ exactly when
\begin{equation}\label{eq:body-count}
 |S|\leq\sum_i\min\{\ell,kp_i(S)\}-\ell
 \qquad(\varnothing\ne S\subseteq J).
\end{equation}
\end{lemma}
\begin{proof}
For necessity, include the whole $A_i$ when $p_i(S)\geq2$,
and include only the one incident pin when $p_i(S)=1$.
Subtracting the mandatory internal edges from the sparsity
bound on this vertex set gives \eqref{eq:body-count}.

For sufficiency, take any vertex set $Y$ and put
$a_i=|Y\cap A_i|$. If $J[Y]\ne\varnothing$, the internal
mandatory contribution is at most $ka_i-\ell$ for $a_i\geq2$,
and zero for $a_i=1$. Apply \eqref{eq:body-count} to $J[Y]$:
\[
 |P_0[Y]|+|J[Y]|
 \leq \sum_{a_i\geq2}(ka_i-\ell)
      +\sum_i\min\{\ell,ka_i\}-\ell
 =k|Y|-\ell.
\]
If $J[Y]$ is empty, mandatory sparsity is enough.
\end{proof}

\begin{corollary}[Constant-rank local oracle]\label{cor:local}
Between two disjoint bodies, $J$ is independent exactly when
\begin{equation}\label{eq:local}
 |J|\leq\ell,\qquad d_J(v)\leq k\ (v\text{ a pin}),\qquad
 |J_{uv}|\leq c\ (u,v\text{ a pin pair}).
\end{equation}
Independence and closure relative to a basis of size at
most $\ell$ are decidable in constant time.
\end{corollary}
\begin{proof}
In \eqref{eq:body-count}, each body's contribution is $k$
for one incident pin and $\ell$ for at least two.
Subsets using at least two pins on both sides have bound
$\ell$; a star has bound $k$; and a single pair has bound
$c$. These are exactly the displayed conditions and their
subset consequences.
\end{proof}

For a fixed body $A_i$, replace all its neighboring bodies
by one independent body, retaining every physical neighbor
pin identity. This gives its \emph{local matroid}, of rank
at most $\ell$, with oracle \eqref{eq:local}.
This operation is used for sampling, not as a global
independence oracle.

\subsection{Logical size and the number of large bodies}

Let $F$ be an independent set of new edges after $P_0$.
Its size is at most $\ell q$, where $q$ is the number
of old bodies. For $q\geq1$, sparsity on the union of all body
vertices and subtraction of their mandatory bases gives even
$|F|\leq\ell(q-1)$; for $q=0$, $F$ is empty.
Throughout the block history
all of $P_0$ is present. Prefixes during the preliminary insertion
of individual mandatory edges are not part of this history;
they are contracted to the single initial timestamp zero.
In a prefix $P_0+F'$, a maximal tight block
$C$ either contains a whole $A_i$ or meets it in at most
one pin, by uncrossing.
Let $b_C$ count the whole old bodies in $C$ and $p_C$
the remaining pins. Define its \emph{logical size} by
$\sigma(C)=b_C+p_C$; this is also the number of old
bodies it meets.

\begin{lemma}\label{lem:logical-size}
For every prefix block,
\begin{equation}\label{eq:logical-bars}
 |F'[C]|=\ell b_C+kp_C-\ell.
\end{equation}
The sum of logical sizes of current blocks is $O(q)$.
For $\sigma(C)\geq2$,
$|F'[C]|\geq(c/2)\sigma(C)$.
For a threshold $s\geq2$, the number $Q$ of persistent
identifiers born for blocks of logical size at least $s$
satisfies
\begin{equation}\label{eq:large-budget}
 Q\leq\frac{2\ell q}{cs}.
\end{equation}
\end{lemma}
\begin{proof}
Subtract the internal mandatory bases of the $b_C$
complete bodies from tightness to obtain
\eqref{eq:logical-bars}.
Let $r$ be the number of current blocks containing whole
old bodies. They partition those $q$ bodies. All current
internal new-edge sets are disjoint, so
\[
 k\sum_{b_C\geq1}p_C+
 \sum_{b_C=0}(kp_C-\ell)
 \leq \ell q-\ell(q-r)=\ell r.
\]
In the second sum, $p_C\geq2$ and
$kp_C-\ell\geq cp_C/2$.
The total pin contribution is therefore $O(q)$.
Also $\ell b_C+kp_C-\ell\geq k\sigma(C)-\ell
\geq c\sigma(C)/2$ when $\sigma(C)\geq2$.

Whenever a new large block contains a previous large
block, inherit one such identifier and retire the others.
Otherwise create a new identifier and reserve the new
block's internal $F$-edges. Birth blocks of different
identifiers have disjoint such edge sets: a shared edge
would force containment in an existing large descendant.
Every birth reserves at least $cs/2$ edges, proving
\eqref{eq:large-budget}.
\end{proof}

The compact prefix history from Section~\ref{sec:sparse}
computes these objects in $q^{1+o(1)}$ time on the actually
registered source graph. Vertices not incident to $F$
remain lazy. A lazy pin can enter a maximal block only
together with its whole old body.

\section{Historical hinges, straddlers, and endpoint coverage}
\label{sec:cover}

\subsection{The total historical hinge count}

\begin{lemma}\label{lem:hinges}
Consider a tracked family of tight sets in an independent
monotone graph. Active sets intersect pairwise in at most
one physical vertex. An identifier follows an increasing
set until it is retired into a containing tracked set.
Let $R$ be the total number of identifiers ever introduced.
The set $T$ of vertices ever belonging to at least two
active tracked sets has
\[
 |T|\leq3\ell R.
\]
After identifying retired body parameters with their
successors, all required point--body attachments and
identifier identifications have an $O(R)$-size history.
\end{lemma}
\begin{proof}
Fix a bounded-indegree orientation of the final independent
graph. Each tracked tight set has at most $\ell$ ports,
because its internal indegree deficit is $\ell$.
A shared vertex is a port of at least one participating
set: disjoint internal edge sets cannot both contain
all its $k$ incoming edges.
Thus at most $\ell R$ vertices are shared at any one time.

A shared vertex ceases to be shared only when its
participating identifiers are identified. At an
identification involving $a\geq2$ identifiers, each
disappearing shared vertex is a port of one of those
sets, so at most $\ell a$ disappear.
There are at most $R$ retirements, and
$\sum a\leq2R$, charging each event to its $a-1$
retired identifiers. The number of first appearances
as a shared vertex is at most the final shared count
plus all disappearances. Hence $|T|\leq3\ell R$.

For the attachment history, track port occurrences
$(v,\text{identifier})$ with $v\in T$.
A retired identifier loses at most $\ell$ such occurrences.
On a surviving increasing identifier, a port disappears
only by becoming saturated. The first saturated owner
of a vertex persists in a growing successor set.
A different saturated owner would require identification:
two simultaneously distinct tight sets cannot both
saturate that vertex. The first saturation is charged
once per vertex, and the other losses to retirements.
The total number of port occurrences is therefore
$O(\ell R+|T|)=O(R)$.

Record an attachment for each new port occurrence.
For a saturated occurrence, an attachment to its first
saturated owner suffices; later parameter identifications
transfer it to successors. This records all necessary
relationships. The prefix history, port lists, and owner
updates extract them in almost-linear preprocessing time.
The argument uses tightness, pairwise one-vertex
intersection, and monotone containment, rather than a
particular size threshold.
\end{proof}

\subsection{Old bodies spanning different large blocks}

An old body is a \emph{straddler} if it is not wholly
contained in a large current block but has at least two
distinct pins in large blocks. Those pins lie in different
large blocks.

\begin{lemma}\label{lem:straddlers}
With $b$ current large blocks and $a$ current straddlers,
\[
 ca\leq\ell(b-1)\qquad(b\geq1).
\]
If $Q$ large identifiers are born in the complete history,
the set $X_{\mathrm h}$ of old bodies ever becoming
straddlers satisfies
\begin{equation}\label{eq:straddler-history}
 |X_{\mathrm h}|\leq(2\ell/c+3\ell)Q.
\end{equation}
\end{lemma}
\begin{proof}
Take the internal bases of the $b$ large blocks and the
$a$ straddling old bodies. They are edge-disjoint subsets
of the independent source graph. Counting their vertex
sets separately overcounts their union by at least $2a$.
Sparsity of the union gives
$2ka\leq\ell(b+a-1)$, which is the first inequality.

A straddler loses that status only when swallowed
entirely by a large block; it can never become a
straddler again. If the union of $d$ old large children
of the swallowing block already contains two pins of
the swallowed body, the snapshot bound charges at most
$\ell(d-1)/c$ such deaths to the children.
The sum of $d-1$ over large identifications is at most $Q$.

Otherwise one of its former large-block pins $z$ lies in
a large block outside the swallowing union. After the
swallowing, $z$ is a shared vertex of two large blocks.
Different swallowed old bodies give different $z$,
because old bodies are disjoint. Lemma~\ref{lem:hinges}
therefore bounds these deaths by $3\ell Q$.
Finally at most $\ell Q/c$ straddlers remain alive.
These three bounds give \eqref{eq:straddler-history}.
The second death case is essential: swallowing a
straddler need not merge two preexisting large blocks.
\end{proof}

\subsection{Logical ports and the potential}

Assign each edge of final $F$ to one of its physical
endpoints so that a pin receives at most $k$ assignments
and all pins of one old body together receive at most
$\ell$. Such an assignment follows from
Lemma~\ref{lem:body-count} and capacitated Hall inequalities.
It is found by an almost-linear sparse flow computation.
Explicitly, use a unit-capacity source arc for each edge of $F$,
arcs from that edge-node to its two pin-nodes, capacity-$k$ arcs
from each pin-node to its old-body node, and capacity-$\ell$ arcs
from body nodes to the sink. For every $S\subseteq F$ the
available endpoint capacity is
$\sum_i\min\{\ell,kp_i(S)\}$, at least $|S|$ by
\eqref{eq:body-count}. Capacitated Hall therefore routes all
edge units. Integral flow gives the asserted assignment on an
$O(q)$-size network of actually registered pins.

The logical objects of a block $C$ are its complete old
bodies, of capacity $\ell$, and its individually intersected
pins, of capacity $k$. By \eqref{eq:logical-bars}, their
internal assignment deficit sums to $\ell$.
At most $\ell$ objects are therefore \emph{logical ports}.
This assignment is distinct from the full-graph orientation
used in Lemma~\ref{lem:hinges}.

Fix $D\geq4\ell$ and
\[
 s=\max\{2,\lfloor D/(4\ell)\rfloor\}.
\]
Let $\mathcal S$ be the inclusion-wise maximal historical
blocks of logical size less than $s$.
Their internal new-edge sets are disjoint, and
Lemma~\ref{lem:logical-size} gives
\[
 \sum_{C\in\mathcal S}\sigma(C)\leq q+2\ell q/c.
\]
Blocks of logical size one are old bodies and contribute
no internal new edges.
Define
\begin{equation}\label{eq:potential}
 p(A_i)=\deg_F(A_i)+
 \ell\sum_{\substack{C\in\mathcal S\\
             A_i\text{ has a logical port object in }C}}
          (\sigma(C)-1).
\end{equation}
If $A_i$ is only partially in $C$, its object is the
single common pin. Since each block has at most $\ell$
logical ports,
\begin{equation}\label{eq:potential-sum}
 \sum_i p(A_i)\leq P_*q,\qquad
 P_*=2\ell+\ell^2(1+2\ell/c).
\end{equation}
Let
\[
 X=\{A_i:p(A_i)\geq D/2\}\cup X_{\mathrm h}.
\]
Equations \eqref{eq:large-budget}--\eqref{eq:potential-sum}
imply $|X|=O(q/D)$ and $Q=O(q/D)$.

\begin{lemma}[Endpoint coverage]\label{lem:cover}
Suppose a chronological sample $H$ has at most $\ell$
edges per old body pair, $F$ is its greedy basis after
$P_0$, and every omitted edge that remains globally open
is spanned by each of $D$ disjoint local layers at both
end bodies. Then every such omitted edge has each endpoint
in an old body from $X$ or in a current large block.
\end{lemma}
\begin{proof}
Let $u\in A_i$, with $A_i\notin X$, and suppose $u$ is
outside all current large blocks. If $A_i$ were wholly
in a large block, this would be impossible.
Since $A_i$ never straddles, at most one of its pins,
say $x$, belongs to large blocks. Thus $u\ne x$.

Consider earlier sample edges at $A_i$.
Edges in $F$ are charged to $\deg_F(A_i)$.
Every other earlier sample edge is spanned by its
$F$-prefix, so both endpoints lie in a prefix block.
A small such block lies in a member $C$ of $\mathcal S$.
The body-pair cap bounds its contribution at $A_i$ by
$\ell(\sigma(C)-1)$.
If its object is a logical port, this is charged to
\eqref{eq:potential}. Among the maximal historical small blocks
in $\mathcal S$, there is at most one in which its object is
saturated: full-body saturation
uses $\ell$ assignments, and two saturated single pins
would use $2k>\ell$ assignments on disjoint internal
edge sets. That one small block contributes at most
$\ell(s-1)\leq D/4$.
Edges spanned in large blocks use pin $x$ at $A_i$.
Hence fewer than $p(A_i)+D/4<3D/4$ earlier sample
edges can use a pin other than $x$.

Every one of the $D$ disjoint local bases spanning the
omitted $uv$ must contain an edge not incident to $x$.
Indeed, a basis using only $x$ has size at most
$k<\ell$, has degree zero at $u$, degree at most
$c<k$ at $v$, and multiplicity zero on $uv$.
The local oracle therefore does not span $uv$.
This requires at least $D$ non-$x$ edges, a contradiction.
The argument is symmetric for the other endpoint.
\end{proof}

\section{The adaptive local sample}
\label{sec:sample}

At a stage, let $q$ be the number of disjoint mandatory
old bodies. Besides remaining original edges, a sorted
stream of $f=O(q)$ known auxiliary edges may arrive.
These \emph{known constraints} will be justified by earlier
root-certificate edges in Section~\ref{sec:interface}.
Within the source graph they are ordinary chronological
edges with distinct element identifiers.

\subsection{Cyclic layers of constant rank}

For each old body keep $D$ local independent sets
$B_1,\ldots,B_D$. An incoming edge goes into the first
local layer not spanning it, if there is one.
At every prefix,
\begin{equation}\label{eq:flat-chain}
 \operatorname{cl}(B_1)\supseteq\cdots
 \supseteq\operatorname{cl}(B_D).
\end{equation}
An edge inserted into a later layer was already spanned
by every earlier one, proving this invariant.
Equal-rank nested flats coincide. Since local rank is
at most $\ell$, group consecutive equal-rank layers
into at most $\ell+1$ intervals, retaining a representative
basis and each layer's actual assigned edges.
Test interval representatives and update the first layer
of the first accepting interval. Only its basis changes;
split the interval and merge equal-rank neighbors.
This costs $O_{k,\ell}(1)$ per processed edge and $O(D)$
initial storage per body.

\subsection{A body-pair filter using only selected edges}

For each old body pair $(a,b)$, maintain $H_{ab}$, the
selected sample edges on that pair. The chronological
reference algorithm on $e=uv$, $u\in A_a$, $v\in A_b$,
is the following.
\begin{enumerate}
 \item If $a=b$, omit $e$ as already spanned by $P_0$.
 If $H_{ab}$ spans $e$ relative to the two mandatory
 bodies, omit it safely.
 \item Otherwise process $e$ in the cyclic local layers
 of both endpoint bodies.
 \item Select $e$ into $H$ and $H_{ab}$ if either endpoint
 layer accepts it, or if it is a known constraint.
 Otherwise leave it outside $H$.
\end{enumerate}
The pair test uses the exact oracle \eqref{eq:local}.
It involves only earlier \emph{selected} edges, rather
than an independently scanned basis of the full input pair.

\begin{lemma}\label{lem:adaptive-sample}
In every prefix, $H_{ab}$ is independent and has size at
most $\ell$, and
\[
 |H|\leq\ell qD+f.
\]
Let $F$ be the chronological greedy basis of $H$ after
$P_0$. Every known constraint is spanned by $P_0$ and
the $F$-prefix including its own source position. A known
constraint omitted by the pair filter is already spanned by the
strictly preceding $F$-prefix. Every omitted original
edge that is not spanned by its $P_0+F$-prefix is spanned
by all $D$ local layers at both endpoint bodies.
\end{lemma}
\begin{proof}
The pair filter inserts only rank-increasing elements
into $H_{ab}$. Each selected original edge uses a local
layer slot at an endpoint; there are $\ell qD$ slots.
At most $f$ additional known edges are selected.
An omitted known edge, or an original edge rejected by
the pair filter, is spanned by earlier $H_{ab}$.
The relevant $F$-prefix is a basis of the $H$-prefix,
so it spans these edges too.
Any remaining omitted original edge was rejected in
every local layer at both endpoints; all those layer
edges belong to $H$. The bases are disjoint within
one endpoint's layer system.
\end{proof}

This proves both hypotheses of Lemma~\ref{lem:cover}.
Known constraints are included in $H$ or its prefix
closure, and $|F|\leq\ell q$ still holds.
No full preliminary scan by body pair is required.

\section{Implicit execution on shared incidence lists}
\label{sec:implicit}

Initially retain only the first $c$ original edges of
each physical vertex pair. The first $c$ copies form
a two-vertex tight set, so this preserves every ordered
greedy prefix. A constant-time initializable direct-address
pair table gives $O(m+n)$ preprocessing \cite{initializable}.
Build the original ordered incidence lists once, using
doubly linked nodes, one at each endpoint of an edge.
Unlinking the two nodes of an edge takes $O(1)$ time.
All stages share these lists; they are never copied.

At a stage, each physical vertex has a view cursor.
Its view starts at the current live list head and can
advance past locally rejected edges without removing
them from the shared list. Those edges remain available
at the next stage. Original selected edges and original
edges with a valid prefix-spanning certificate are
removed permanently.
Known auxiliary edges use separate small ordered lists.
Stage membership and activation predicates are fixed;
their safety is proved in Section~\ref{sec:interface}.

\subsection{Finding the next locally eligible edge}

\begin{lemma}\label{lem:successor}
For a nonfull final local layer $B_D$ and a fixed own
pin $u$, its next locally eligible original incidence
is found with constant local skipping work, apart from
edges deleted permanently. A saturated own pin is
disabled as a whole. Changes in $B_D$ cause at most
$\ell$ full point-head rebuilds per body.
\end{lemma}
\begin{proof}
If $|B_D|=\ell$, disable local search for the entire
body. Otherwise, an edge is locally rejected exactly
because its own pin has degree $k$, its neighbor pin
has degree $k$, or its physical pair has multiplicity $c$.
There is at most one saturated own pin because
$|B_D|<\ell<2k$; disable its head in constant time.
For another own pin $u$, at most $\ell$ neighbor
identifiers are forbidden. Each has at most $c$ original
incidences with $u$ in the entire root input.
Thus at most $c\ell$ locally forbidden original records
are skipped before the next eligible record.
In fact this also bounds their total number at one own pin
throughout a stage before that pin or its body is disabled:
$B_D$ only gains edges, its union of neighbor identities has size
at most $\ell$, and a view cursor never moves backwards.

The last layer grows at most $\ell$ times, and its
flat is unchanged between growth events. On each growth,
recompute all point heads in that body, continuing from
the current view cursors rather than resetting them.
The separate known-edge lists contain $O(f)$ records;
their total cursor advancement is $O(f)$.
\end{proof}

Keep an inner heap of eligible point heads for each
body and an outer heap of the body minima. Compare
the outer minimum with the next known constraint.
Process the earliest edge that is locally eligible
at some endpoint and still open in its body pair,
or the next known constraint.
When an endpoint's last layer grows, rebuild that
body's point heads. Otherwise advance only the affected
endpoint's head.

Pair closure may make a queued head stale. Repair it
lazily: remove a whole chain of currently pair-spanned
incidences, perform all their constant-time unlinks,
and insert only the next valid head into the heap.
There is no individual heap update for each deleted
edge. If the other endpoint's view cursor points at
the removed node, advance it in constant time.
A queued copy at that endpoint may remain stale until
its own later repair.
Keep at most one queued record per point and keep its key
separate from the live view cursor. A stale queued node is
detected by its live flag; repairing it starts from the already
advanced cursor. Thus deletion does not require updates of
other points' heap entries, or traversal of chains of dead nodes.

The order of tests during head generation is essential.
At either endpoint, check the static safety predicates of
the entire edge (including both activation thresholds and both
body identifiers), then local closure,
and use the current $H_{ab}$ test only on a locally
eligible record. An earlier locally skipped record is
never deleted using a later sample edge as a witness.
Local closure only grows, so that record cannot become
eligible again in this stage. Every still-live eligible
record deleted by a pair test occurs after the already
processed sample edges. Pair deletion therefore always
has a chronological certificate.

\Needspace{8\baselineskip}
\subsection{Saturations affecting many heads}

\begin{lemma}\label{lem:saturation-events}
While an old body $A_a$ has a nonfull final local layer,
at most $\ell D+1$ different pair bases $H_{ab}$ can
reach rank $\ell$. At most $\ell D+\ell+1$ different
neighbor pins $y$ can reach degree $k$ in their pair
base with $A_a$.
\end{lemma}
\begin{proof}
Take the final still-nonfull basis $B$ of the last
layer during this active interval: immediately before its first
full-rank insertion, or at the end of the stage if it never fills.
A full pair basis has rank $\ell$ in the local matroid
as well: on that one pair the local and pair oracles
coincide. It contains an edge outside $\operatorname{cl}(B)$.
Since layer flats only grow, that edge was accepted
in a local layer of $A_a$ on arrival.
Different body pairs give distinct witnesses.
There are at most $\ell D$ local acceptances; the
final edge that disables the body adds at most one event.

A degree-$k$ pair star at a neighbor pin $y$ either
contains an edge outside $\operatorname{cl}(B)$,
giving another local acceptance, or is wholly spanned
by $B$. Stars for different $y$ have disjoint edges.
In the latter case $y$ occurs in $B$.
Indeed, all $k$ star edges cannot use the single
saturated own pin, because each physical pair has
at most $c<k$ selected copies.
A spanned star edge from another own pin therefore
requires saturation of $y$ in $B$ or saturation of
its physical pair in $B$. Both make $y$ a neighbor
endpoint of a basis edge.
There are at most $\ell$ such $y$. Again the disabling
edge adds at most one event.
\end{proof}

\begin{lemma}[Implicit sampling cost]\label{lem:heap-cost}
Let $p$ be the number of physical vertices handled at a
stage, and let $d$ be the number of original edges
permanently unlinked at that stage. The implicit procedure
produces exactly the reference sample and local layers,
in time
\begin{equation}\label{eq:heap-cost}
 O\bigl(d+(Dp+qD+f)\log(n+2)+f\bigr).
\end{equation}
For the fixed $D$ used below, the coefficient of $d$
does not depend on the number of stages.
\end{lemma}
\begin{proof}
Pair closure changes only on four kinds of saturation:
total rank $\ell$, own-pin degree $k$, neighbor-pin
degree $k$, and physical-pair multiplicity $c$.
The first and third kinds may affect all $p_a$ point
heads of a body, but Lemma~\ref{lem:saturation-events}
bounds them by $O(D)$ events while that body is active.
The other two kinds affect only one own point head,
and there are $O(|H|)$ such events.

A repaired head is checked against \emph{all current}
pair closure conditions. Another invalidation therefore
needs a later saturation event or selection of that
edge. This remains true if another endpoint's repair
deletes the queued edge: an already valid head can be
deleted only after a new pair-spanning event, or after
selection. Static predicates are checked at both
endpoint head generations and cannot later invalidate
an otherwise valid head.
Thus there are $O(Dp+|H|)$ head repairs and
$O(\ell p)$ rebuilding point visits.
Binary heaps give the logarithmic term, since
$|H|\leq\ell qD+f$.
More explicitly, charge a stale point head to the first event
since its last valid installation that invalidated that record.
A rank or neighbor-pin saturation is charged at most once per
point before that point's next repair, and there are $O(D)$ such
events per active body. An own-pin or physical-pair saturation
affects just its own point in each endpoint body. Selection of
an original edge is charged to at most its two queued copies.
Rebuilding a body's heads on growth of $B_D$ accounts separately
for at most $\ell$ visits to each of its points. Deleting several
records during one repair does not create an event or a heap
operation for each record.

Locally skipped original incidences are bounded by
Lemma~\ref{lem:successor}. A long safe-deletion chain
charges each edge only its constant-time unlink;
it has only one final heap update.
Known-list cursor work sums to $O(f)$.

The heaps and known stream always expose the earliest
edge eligible at an endpoint and open in its pair.
An original edge skipped at both endpoints is already
rejected by all layers there. Monotone flats cannot
make it eligible later.
An earlier eligible edge therefore cannot be hidden
behind a later processed edge.
Induction over selected edges gives exactly the
reference sample and exactly the reference local
layer contents.
\end{proof}

\section{The fixed first-binding interface}
\label{sec:interface}

Sampling and the sparse optimizer give the chronological
basis $F$ of $P_0+H$. From it choose all current large
blocks and the fixed exceptional old bodies $X$.
Retire an exceptional body's identifier when that
body is wholly absorbed by a large block.
The active tracked tight sets intersect pairwise in
at most one pin. Their total number of persistent
identifiers is
\[
 R\leq Q+|X|=O(q/D).
\]
Lemma~\ref{lem:hinges} applies to this family, including
the exceptional bodies as identifiers present initially.
It gives $O(R)$ known relationships and a set $T$ of
historically shared physical pins with $|T|\leq3\ell R$.

\begin{lemma}\label{lem:interface}
In $q^{1+o(1)}$ preprocessing time and $O(n)$ additional
original-vertex work, the source history can be turned
into the next stage with:
\begin{enumerate}
 \item $q'\leq C_*q/D$ disjoint mandatory tight frames;
 \item $O(q')$ known chronological auxiliary edges;
 \item one fixed body identifier $\lambda'(v)$ and one
 activation threshold $\tau'(v)$ per original vertex; an always
 inactive vertex uses $\lambda'(v)=\botval$.
\end{enumerate}
Here $C_*$ depends only on $k,\ell$.
Closure in the new model is safe in the corresponding
original weight prefix. An unselected original edge
that fails its endpoint activation test is already
spanned by earlier root-certificate edges.
\end{lemma}
\begin{proof}
\emph{Disjoint initial frames.}
Make one fresh tight frame for every tracked identifier.
For every $z\in T$, also make a distinct fresh tight
frame with exactly one used physical pin, namely $z$,
and its $k$ private attachments.
The pin retains $k$ free coordinates, because $U_z$
has row rank $k$; the extra $\ell-k$ frame parameters
impose no original constraint.
Initial frames are disjoint and their internal bases
are independent. There are
\begin{equation}\label{eq:child-count}
 q'=R+|T|\leq(1+3\ell)R
\end{equation}
of them.

\emph{Known constraints.}
An actual attachment of a shared pin $z$ to tracked
frame $A$ is represented by $k$ edges from $z$ to
distinct core vertices of $A$. By \eqref{eq:motions},
these impose $\xi_z=U_z\alpha_A$.
To identify a retired frame $B$ with frame $A$,
use $\ell$ diagonal edges between their cores,
with at most $c$ copies per corresponding core pair.
The width \eqref{eq:frame-width} is sufficient.
Their degrees are at most $c<k$, their pair
multiplicities at most $c$, and their number is $\ell$.
By \eqref{eq:local}, they are independent and make
the two-frame union tight in the contraction by the two mandatory
frames. Their constraints annihilate all motions with equal
parameters. There are $\ell$ independent constraints on the
$2\ell$ parameters, so their nullspace is exactly that
$\ell$-dimensional equal-parameter space. They impose exactly
$\alpha_A=\alpha_B$.
No fresh merger pins are needed.
Insert these known edges at the source $F$-prefix
when the relationship becomes true.
The compressed history has $O(R)=O(q')$ such edges.

\emph{Fixed bindings and times.}
For an original pin $v\notin T$, bind it to the
\emph{first} tracked identifier that contains it.
If no such identifier exists, leave it unbound and inactive.
This identifier is fixed for the whole next stage.
Later retirement is handled by parameter equality,
so no endpoint renaming is necessary.
If $v$ could belong to two distinct active tracked
sets, the complete offline history would place it
in $T$; then its separate free frame and known
attachments retain the same physical coordinates.
For $v\in T$, use that separate free frame as its
fixed initial body.

Let $b(v)$ be the first tracked membership time;
put $b(v)=+\infty$ when there is no membership.
In the implementation $+\infty$ is the sentinel $m+1$, which
is larger than every original query index.
Compose the new threshold with the old one by
\begin{equation}\label{eq:activation}
 \tau'(v)=\max\{\tau(v),b(v)\}.
\end{equation}
Frames that will be born later may already exist
initially: no earlier usable original edge involves
them, and their known relationships arrive only at
the justified prefix.
Coordinates of an original vertex bound to such a
frame need not match the root coordinates before
its activation time.

\emph{The chronological extension invariant.}
For each root weight prefix, every motion satisfying
the selected original certificate has a lift to a motion
of the current source model, with identical coordinates
at all original vertices whose thresholds have passed.
Each source model has its own coordinate space. It uses the same
$U_v$ and original-edge coefficients for an original physical
identity $v$, but agreement of its coordinate vector with the
root vector is imposed only after activation. Before activation
the source coordinate may be chosen independently of the root
coordinate. Thus ``lift'' does not assert simultaneous satisfaction
of all stages' mandatory rows in one root coordinate space.
The source rows are its mandatory initial bases,
the known constraints arrived so far, and its own
selected original sample rows so far.
Original rows from ancestor stages are not copied
into the new source model; the root motion already
satisfies them.

To establish the invariant at the interface, first take
a lift of the root motion to the parent model.
Its $F$-prefix is a basis of its $H$-prefix.
On each tracked tight set, the motion thus has
form \eqref{eq:motions}; choose that $\alpha_A$
for its new frame. For a persistent identifier its birth set
already determines the parameter uniquely. Its growing sets
and containing successor sets already present at this prefix
have the same parameter,
because they contain that birth set. Thus these choices are
consistent with the arrived history. An identifier not yet born
has no arrived relationship and its parameter can be chosen freely.
Preserve the physical coordinates of shared pins.
Attachments and parameter identifications hold
at exactly their justified parent times.
Unused future frame coordinates can be completed
freely. If a root vertex is active in the child,
it was active in the parent and already belongs
to its tracked set, so its coordinates agree.
Every sample edge selected subsequently at the
child stage has active endpoints on arrival.
Thresholds are fixed within a stage, so those rows
continue to be evaluated on root coordinates in
all later prefixes and are satisfied by the root
certificate motion.
This proves the invariant and hence the safety
of source closure.

\emph{Activation discards.}
Consider an original edge outside $H$ that passed
its old thresholds but has an endpoint not yet in
any tracked set. If it were open in the parent
$P_0+F$-prefix, Lemma~\ref{lem:cover} would place
both endpoints in the tracked family.
It is therefore already parent-prefix-spanned.
The old invariant still identifies its coordinates
with the root coordinates, so this discard is safe.
Failure of an earlier threshold was already justified
at an earlier interface.

\emph{Computing vertex metadata.}
Only $O(q)$ physical pins occur in actual $F$-edges.
A never-used lazy pin cannot partially enter a
new maximal block: without new incident edges,
its unique maximal block contains its whole old
body. Such pins therefore share their old body's
first-membership event.
Registered-pin exceptions and $T$ are obtained
from the $O(q)$-size actual source history.
For each original vertex, its old identifier and
optional registered-pin record determine its new
identifier and threshold. One $O(n)$ vertex pass
computes all metadata without reading original
edges. The sparse history preparation costs
$q^{1+o(1)}$.
\end{proof}

\begin{remark}
The child need not inherit the closure of every small
parent block. Such an edge may be open in the weaker
child model and may be selected into a later sample.
This is allowed: each stage selects at most
$\ell qD+f$ original and known edges, and final
optimization is performed in the original matroid.
The argument uses safe discards, not equivalence of
all parent and child closure answers.
\end{remark}

\section{One fixed algorithm and its total cost}
\label{sec:complete}

\subsection{Fixing the constants once}

Use the constants
\begin{align*}
 P_*&=2\ell+\ell^2(1+2\ell/c),\\
 A_*&=16\ell^2/c,\\
 C_R&=A_*+2P_*+(2\ell/c+3\ell)A_*,\\
 C_*&=(1+3\ell)C_R,
\end{align*}
and fix
\begin{equation}\label{eq:fixed-D}
 D=\left\lceil\max\{4\ell,4C_*\}\right\rceil.
\end{equation}
The threshold $s=\max\{2,\lfloor D/(4\ell)\rfloor\}$
satisfies $s\geq D/(8\ell)$.
Thus $Q\leq A_*q/D$.
The potential exceptions contribute at most $2P_*q/D$,
and the historical straddlers at most
$(2\ell/c+3\ell)Q$.
It follows that $R\leq C_Rq/D$, and
Lemma~\ref{lem:interface} gives
\begin{equation}\label{eq:geometric}
 q'\leq C_*q/D\leq q/4.
\end{equation}
The integer $D$ depends only on $k,\ell$, not on $n$,
$m$, recursion depth, or an exponent-error parameter.

\subsection{The initial stage}

Give each original vertex $v$ a disjoint constant-size
tight frame with $v$ as its one used physical pin,
attached privately by $k$ edges.
Start with $q_0=n$, $\tau_0(v)=0$, and no known
auxiliary edge.
This introduces no restriction on original coordinates:
one used pin retains $k$ free coordinates in its
$\ell$-parameter body.
Equivalently, in \eqref{eq:body-count} every touched
body contributes $k$, so the induced edge matroid
is exactly the original $(k,\ell)$-sparsity matroid.
Initial auxiliary rows are not returned as original edges.

\begin{algorithm}[htbp]
\caption{Uniform deterministic original-edge certificate}
\label{alg:complete}
\textbf{Input:} Ordered original edges; fixed $k<\ell<2k$.\par
\textbf{Output:} $\mathcal C$ with $|\mathcal C|=O(n)$ and
$\Bj\ell{\mathcal C}=\Bj\ell E$.
\begin{enumerate}[label=\arabic*:,leftmargin=2em,itemsep=2pt,topsep=4pt]
\item Fix $D$ by \eqref{eq:fixed-D}.
\item Keep only the first $c$ copies of every physical pair.
\item Build the shared ordered incidence lists once.
\item $\mathcal C\gets\varnothing$; initialize the $n$-body model.
\item \textbf{While} the current body count $q$ exceeds $D$:
\begin{enumerate}[label=\alph*.,leftmargin=1.5em,itemsep=2pt,topsep=2pt]
 \item Generate $H$ by the implicit adaptive-pair sampler.
 \item Add every selected original edge to $\mathcal C$.
 \item Permanently unlink selected or safely spanned original edges.
 \item On the explicit sparse graph, compute the chronological
 basis $F$ after the mandatory $P_0$.
 \item Compute logical block history, $X$, tracked identifiers,
 and shared pins $T$.
 \item Build the child model and fixed vertex metadata by
 Lemma~\ref{lem:interface}.
\end{enumerate}
\item In one final pass over the original edge array, process
 still-live usable edges in the constant-size source model,
 merging the known-constraint stream chronologically.
\item Add its rank-increasing original edges to $\mathcal C$.
\item \textbf{Return} $\mathcal C$.
\end{enumerate}
\end{algorithm}

At stage $j$, the metadata are fixed. An original
$e_i=uv$ is usable precisely when
$i\geq\tau_j(u),\tau_j(v)$.
All ancestor activation conditions have already
been folded into these two thresholds.
Only two endpoint lookups are performed; the edge
does not query a sequence of ancestor closure indices.
Within-body edges are spanned by the mandatory
model. Static usability and within-body tests are
applied during implicit head generation.
Edges at vertices whose views never become active
need not be discovered until the final array pass.

Known constraints carry root-prefix timestamps. Such a constraint
may be justified by several selected original edges; its timestamp
identifies the last possible supporting original position, rather
than claiming that one original edge alone implies the constraint.
Every positive such timestamp is the index of an original edge
selected at an ancestor or earlier stage and already removed
from the live lists. The same is true of every positive first-
membership and activation timestamp.
Constraints with that timestamp are ordered
after that original edge and before the next
original edge. An inherited known-constraint time
has the same property transitively.
Thus no still-live original query at the same
root index precedes a constraint it uses.
Auxiliary edge elements still have distinct identifiers.
For a still-live original $e_i$, passing a finite positive
activation threshold therefore means that threshold is strictly
less than $i$. All source conditions used to discard $e_i$ are
valid in the root prefix strictly preceding $e_i$, not merely in
a prefix that includes the queried edge. Timestamp zero refers
to the unrestricted initial model. Multiple auxiliary edges at
one timestamp have a fixed internal order; original queries at
that timestamp have already been removed.

\subsection{The constant-body terminal stage}

When $q\leq D$, the source has constant total body
parameter dimension. Use the incremental component
matrix of the component pebble-game algorithm
\cite{quadratic}, on a preallocated constant-size
vertex universe containing all possible successful
registrations.
Pre-register the endpoints of the $O(q)$ known rows, as well as
the cores. Reserve at most $2\ell q$ additional original-pin slots.
There are at most $\ell q$ independent new source
rows and $O(q)$ known rows, and consequently only
$O(q)$ actually registered query endpoints.

If $q=0$, all remaining original vertices are inactive and the
final pass only performs the activation test. For $q>0$, apply
both-endpoint activation and within-body tests first.
An unsuccessful virtual-pin query requires no
registration:
\begin{enumerate}
 \item For two registered endpoints, read their
 common-component bit.
 \item For a registered $u$ and an unregistered
 pin attached to frame $B$, test the closure
 bits from $u$ to the $k$ distinguished core
 vertices of $B$. They are all true exactly
 when $u$ belongs to the maximal block
 containing $B$.
 \item For two unregistered pins of frames
 $A,B$, test all cross-pair core closure bits.
 They are all true exactly when the two
 frames lie in one maximal block.
\end{enumerate}
The frame size is constant.
For the second equivalence, the $k$ tested edge rows, together
with a tight basis of $B$, form a tight graph on $B\cup\{u\}$
(membership $u\in B$ is immediate). If all tested rows are
spanned, Lemma~\ref{lem:spanned-tight} puts $u$ and $B$ in one
block. The converse is Lemma~\ref{lem:uncross}.
For the third equivalence, all true cross-pair bits span the
$\ell$ diagonal rows used to join $A$ and $B$; those rows and
the two frame bases form a tight graph. Apply
Lemma~\ref{lem:spanned-tight} again. Finally a pin with no new
incident edge belongs only to the maximal block containing its
frame, by the fresh-pin argument in Section~\ref{sec:bodies}.
These statements also show that omitting unsuccessful
registrations leaves all answers on registered vertices unchanged.
On a successful original query, register its
missing endpoints with private attachments
and insert the original edge.
The bounded total number of insertions costs
$O(1)$ for fixed $k,\ell,D$, and each remaining
original query has constant cost.
The final pass therefore costs $O(m+n)$.

\subsection{Correctness and the uniform additive bound}

\begin{proof}[Proof of Theorem~\ref{thm:main}]
For $\ell\leq k$, use Proposition~\ref{prop:lower}.
For the upper range, construct $\mathcal C$ by
Algorithm~\ref{alg:complete} and then run
$\Sparse(\ell,\mathcal C)$ in the original matroid.

\emph{Correctness.}
The initial physical-pair truncation preserves
ordered greedy prefixes, since each later copy
is spanned by the first $c$ copies.
On the retained input, every later original
discard has a certificate using only earlier
selected original edges in its own weight prefix.
Pair-filter discards use earlier $H_{ab}$.
Activation discards and within-body discards
are justified by Lemma~\ref{lem:interface}.
The motion-extension invariant transfers every
source dependence to the original matroid.
Known auxiliary conditions themselves are
supported by earlier original samples, transitively;
they are not included in the original output.

A weaker child model may select an original
edge that was already dependent in the parent.
Keeping such an extra candidate is harmless:
all samples are retained in $\mathcal C$, and
final optimization uses the original matroid.
At the terminal stage, the exact source basis
and its known constraints justify all remaining
discards in the same way.
Thus every omitted original edge is spanned
by the earlier $\mathcal C$-prefix.
This preserves every prefix rank and hence
the ordered greedy output. The final sparse
optimization returns the required maximum-weight
$(k,\ell)$-sparse subgraph.
Initial truncation witnesses also transfer to
the final certificate by this prefix-preservation
property.

\emph{Original input work.}
Reading and truncating the input and building
the incidence lists costs $O(m+n)$.
Let $d_j$ count original unlinks at stage $j$.
Each edge is unlinked at most once, so
\[
 \sum_j d_j\leq m.
\]
Its two list modifications take constant time.
Lemma~\ref{lem:heap-cost} does not attach a
separate heap update to every deleted edge.
The single final array pass adds $O(m)$.
Original records skipped locally without
deletion are charged to the stage's point-head
bound, rather than to a repeated complete
edge scan.
The direct original-edge work is therefore
$C_{k,\ell}m$, with a coefficient independent
of recursion depth.

\emph{Point and source work.}
By \eqref{eq:geometric},
\[
 \sum_jq_j=O(n),\qquad
 \text{number of stages}=O(\log(n+2)).
\]
Known-edge counts satisfy $f_j=O(q_j)$,
so $\sum_jf_j=O(n)$.
The explicit graph for a stage sample has
$O(Dq_j+f_j)$ vertices and edges, including
its constant frames and registered pins.
After computing $F$, unused physical pins
can again be treated lazily.
Since $D$ is fixed, all sparse optimizations
and prefix-history constructions together cost
\[
 \sum_j(Dq_j+f_j)^{1+o(1)}
 =n^{1+o(1)}.
\]

A stage handles at most the $n$ original
points and $O(q_j)$ known-edge and frame
points. Thus $p_j\leq n+O(q_j)$ and
$\sum_jp_j=O(n\log(n+2))$.
Vertex-metadata passes, head rebuilds, and
the non-input part of \eqref{eq:heap-cost}
sum to $O(n\log^2(n+2))$.
All freshly created auxiliary points over
all stages number $O(\sum_jq_j)=O(n)$.

Finally,
\[
 |\mathcal C|
 \leq\sum_j(\ell Dq_j+f_j)+O(1)=O(n).
\]
Its final sparse optimization takes
$n^{1+o(1)}$ time.
Combining the bounds gives
\[
 O\!\left(m+n\log^2(n+2)
       +\sum_jq_j^{1+o(1)}\right)
 =O(m+n^{1+o(1)}).
\]
All heaps, local oracles, initializable tables,
flows, decompositions, and history computations
are deterministic. The constants were fixed
once in \eqref{eq:fixed-D}, establishing the
claimed uniform quantifier order.
\end{proof}

\begin{remark}
The generic matrices and motions occur only
in the correctness proof. No random coordinates
or polynomial identity tests are used.
Body-pair tables may use quadratic address space
with constant-time initialization. The claimed
bound is the uniform additive running time.
\end{remark}

\section{Exact checks and their scope}
\label{sec:checks}

The original manuscript reports the following exact combinatorial
checks from its development. Their scripts and logs were not supplied
with the manuscript, so these original counts have not been
independently reproduced in this review.
\begin{itemize}
 \item For fixed body partitions, $160$ implicit-local
 sampling tests, with $72\,533$ input edges in total,
 produced exactly the same local layers and samples
 as direct chronological execution.
 \item The adaptive selected-pair filter was checked
 on another $160$ tests with $120\,414$ input edges.
 All layer contents and selected edges agreed.
 These runs had $43\,497$ safe deletions and only
 $3516$ lazy heap-head repairs. The saturation-event
 bounds were checked for every active body.
 A separate heavy-pair example had $3765$ input
 edges, $33$ selected edges, and $3732$ safe deletions.
 \item Endpoint coverage was checked on four exact
 source graphs with $2014$ input edges.
 A further structured $601$-edge example had a
 genuinely open omitted edge: the sample contained
 the first $600$ edges, and the last edge's endpoints
 were covered by two different large blocks.
 \item Fixed first bindings, separate free shared-pin
 frames, and timed diagonal identifications were
 checked on $24$ exact histories.
 There were $4046$ helper queries without an external
 parent-closure guard, including $952$ edges already
 spanned in the parent. Another $172$ edges were
 safely omitted by activation times before their
 future frames were born.
 Every resulting original-edge certificate gave
 the same ordered greedy optimum as the full input.
\end{itemize}

The October 11, 2026 review additionally supplies reproducible
Python scripts using integer arithmetic, exhaustive vertex-subset
tests, and integral minimum cuts. With fixed seeds $20261011$
and $20261012$, all of the following independent checks passed:
\begin{itemize}
 \item $2158$ weak/strong circuit and greedy-compatibility instances,
 including both edge orders on every four-vertex multigraph with
 zero, one, or two copies per pair for $(k,j)=(2,3)$;
 $1429$ prefix histories and $83\,992$ closure queries.
 \item $1800$ disjoint-body independence queries against exhaustive
 physical-vertex sparsity; $250$ adaptive-sampling instances with
 $43\,093$ input edges and $1040$ separate pair-discard safety checks.
 \item $250$ implicit/reference sampling comparisons, with
 $43\,686$ original edges and $2189$ timed known rows.
 Samples and every local layer agreed. The runs checked
 $26\,460$ safe unlinks and $7665$ stale-head repairs.
 The test uses a flat heap of point heads, which exposes the same
 minimum as the nested heaps; it does not measure their runtime.
 \item $200$ reverse rooted-cut histories with $1485$ states,
 compared with exhaustive cut enumeration.
 \item $250$ lower-range union-certificate instances, including
 $188$ cases with $\ell=0$ or $\ell=k$. The final-present-type
 shortcut chose the same layer as direct execution, and every
 certificate preserved the full ordered greedy output.
 \item $100$ covered tracked-frame interface histories, involving
 $278$ identifiers, $132$ historically shared pins,
 $1096$ known-row queries and $555$ remaining original queries.
 All $441$ discards, including $85$ activation discards, were
 safe in the strict parent prefix, and every certificate preserved
 the ordered parent greedy output.
 These tests augment the exceptional family to enforce qualitative
 endpoint coverage; they do not test its asymptotic size bound.
 \item $3796$ terminal virtual-pin closure queries, compared with
 the exact body-count oracle after explicit virtual-pin inclusion.
\end{itemize}

The scripts and their machine-readable result files accompany the
reviewed source. These finite checks exercise specific lemmas and
interfaces. They are not a full implementation, formal verification,
or an empirical verification of the asymptotic running time.
The running-time claim rests on the amortized proof of
Theorem~\ref{thm:main} and the cited algorithmic primitives.

\begin{thebibliography}{99}
\raggedright

\bibitem{initializable}
Torben Hagerup and Frank Kammer.
\emph{On-the-Fly Array Initialization in Less Space}.
ISAAC 2017, LIPIcs 92, 44:1--44:12.
\url{https://doi.org/10.4230/LIPIcs.ISAAC.2017.44}.

\bibitem{natural}
Ileana Streinu and Louis Theran.
\emph{Natural Realizations of Sparsity Matroids}.
Ars Mathematica Contemporanea 4(1), 141--151, 2011.
See Section 3 for the generic $U$ construction.
\url{https://arxiv.org/abs/0711.3013}.

\bibitem{mcf}
Jan van den Brand, Li Chen, Rasmus Kyng, Yang P. Liu,
Richard Peng, Maximilian Probst Gutenberg, Sushant Sachdeva,
and Aaron Sidford.
\emph{A Deterministic Almost-Linear Time Algorithm
for Minimum-Cost Flow}.
FOCS 2023; arXiv:2309.16629.
\url{https://arxiv.org/abs/2309.16629}.

\bibitem{forest}
Pavel Arkhipov and Vladimir Kolmogorov.
\emph{Faster algorithms for packing forests in graphs
and related problems}.
arXiv:2409.20314, version 3, 2026.
\url{https://arxiv.org/abs/2409.20314}.

\bibitem{recognition}
Bence Deák and Péter Madarasi.
\emph{Asymptotically Faster Algorithms for Recognizing
$(k,\ell)$-Sparse Graphs}.
arXiv:2604.13025, version 2, 2026.
\url{https://arxiv.org/abs/2604.13025}.

\bibitem{arborescences}
Anand Bhalgat, Ramesh Hariharan, Telikepalli Kavitha,
and Debmalya Panigrahi.
\emph{Fast Edge Splitting and Edmonds' Arborescence
Construction for Unweighted Graphs}.
SODA 2008, pp.\ 455--464.
\url{https://users.cs.duke.edu/~debmalya/papers/soda08-splitting.pdf}.

\bibitem{arbo-study}
Loukas Georgiadis, Dionysios Kefallinos, Anna Mpanti,
and Stavros D. Nikolopoulos.
\emph{An Experimental Study of Algorithms for Packing
Arborescences}.
SEA 2022, LIPIcs 233, 14:1--14:16.
\url{https://doi.org/10.4230/LIPIcs.SEA.2022.14}.

\bibitem{minset}
Harold N. Gabow.
\emph{The minset-poset approach to representations
of graph connectivity}.
ACM Transactions on Algorithms 12(2), 24:1--24:73, 2016.
\url{https://doi.org/10.1145/2764909}.

\bibitem{minset-2026}
Loukas Georgiadis, Evangelos Kipouridis, Evangelos Kosinas,
Charis Papadopoulos, and Nikos Parotsidis.
\emph{Computing the $(k+2)$-Edge-Connected Components in
$k$-Edge-Connected Digraphs in Subquadratic Time}.
ICALP 2026, LIPIcs 374, 95:1--95:23.
See Proposition 27 for the minset-poset construction bound.
\url{https://arxiv.org/abs/2604.27474}.

\bibitem{static-union}
Harold N. Gabow and Robert Endre Tarjan.
\emph{A linear-time algorithm for a special case
of disjoint set union}.
Journal of Computer and System Sciences 30(2),
209--221, 1985.
\url{https://doi.org/10.1016/0022-0000(85)90014-5}.

\bibitem{quadratic}
Bence Deák and Péter Madarasi.
\emph{Quadratic-Time Algorithm for the Maximum-Weight
$(k,\ell)$-Sparse Subgraph Problem}.
arXiv:2511.20882, 2025.
\url{https://arxiv.org/abs/2511.20882}.

\end{thebibliography}
\end{document}
